\section{The matching endpoint lower construction}
\label{sec:ext-weights}
This appendix reproduces the Phase~8 endpoint construction~\cite{huang8},
with its full kernel proof and bounded-loss normalization. Combined with
Theorem~\ref{thm:weighted-sharp}, it proves the matching sharp lower order.
The source weights and posterior gaps are both allowed to vary with degree.

\subsection{An endpoint-localized polynomial curvature}
The only approximation ingredient needed is elementary external growth of
Chebyshev polynomials. It is a classical polynomial localization mechanism;
the identity $T_n(\cosh v)=\cosh(nv)$ is recorded in
\cite[Eq.~18.5.1]{dlmf}. Squared-Chebyshev endpoint needles and their exponential tail mechanisms
also appear explicitly in Slot--Laurent~\cite[Section~4.1]{slot2019};
see also~\cite{slot2020}. None of this polynomial technology is claimed
as original here. The kernel estimates below are supplied in full.

For an integer $n$ sufficiently large, put
\begin{equation}
 d=d_n=\left(\frac{4\log n}{n}\right)^2,\qquad
 p_n(t)=T_n\!\left(\frac{2t-1-d}{1-d}\right)^2,\qquad
 Z_n=\int_0^1p_n(t)\rd t,\qquad
 \kappa_n(t)=\frac{p_n(t)}{Z_n}.
 \label{eq:ext-endpointneedle}
\end{equation}
Here $\kappa_n$ denotes its polynomial on the whole real line, although its
normalization and probability-mass estimates refer only to $[0,1]$.
Its degree is $2n$ and it is nonnegative everywhere. For $t\in[d,1]$ the
argument of $T_n$ belongs to $[-1,1]$, so $p_n(t)\le1$. On $[0,d/2]$ the
absolute value of that argument is at least $1+d/(1-d)\ge1+d$.
For $0\le v\le1$, $\operatorname{arcosh}(1+v)\ge\sqrt v$: for example,
$\cosh(\sqrt v)\le1+v$ follows by its power series. Using the parity of
$T_n$ and the hyperbolic identity therefore gives, for $d\le1/4$,
\[
 p_n(t)\ge\tfrac14e^{2n\sqrt d}\quad(0\le t\le d/2),\qquad
 Z_n\ge\frac d8 e^{2n\sqrt d}.
\]
Consequently
\begin{equation}
 \int_d^1\kappa_n(t)\rd t\le\tau_n,
 \qquad
 \tau_n:=8d^{-1}e^{-2n\sqrt d}
       =\frac{1}{2n^6(\log n)^2}=o(d^3).
 \label{eq:ext-tail}
\end{equation}
The normalized maximum also satisfies
\begin{equation}
 \max_{[0,1]}\kappa_n\le32n^2.
 \label{eq:ext-max}
\end{equation}
For completeness, the rescaled classical Markov inequality
\cite{shadrin} gives $\|\kappa_n'\|_\infty\le2(2n)^2M$ when
$M=\|\kappa_n\|_\infty$. From a maximizer, one direction remains in $[0,1]$
for length $1/(4(2n)^2)$, throughout which $\kappa_n\ge M/2$.
Integrating and using $\int_0^1\kappa_n=1$ yields $M\le8(2n)^2$.

Define
\begin{equation}
 g_n(t)=\kappa_n(t)+\kappa_n(1-t)+d^2,
 \qquad \ph_n''=g_n,
 \qquad \ph_n'(1/2)=\ph_n(1/2)=0.
 \label{eq:ext-curvature}
\end{equation}
This curvature is strictly positive on the whole real line and symmetric
about $1/2$. Its degree is exactly $r=2n$: both squared-polynomial terms
have positive leading coefficient of the same even degree. The generator
is a symmetric strictly convex polynomial of degree $2n+2$.

\subsection{All deterministic hierarchies are expensive}
Take the sorted source locations and probabilities
\begin{equation}
 (x_A,x_B,x_C,x_D)=(0,2d,1-2d,1),\qquad
 (p_A,p_B,p_C,p_D)=(a,b,b,a),\qquad
 b=d^2,\quad a=\tfrac12-b.
 \label{eq:ext-source}
\end{equation}
For all sufficiently large $n$, these are distinct and all weights are
positive. Every distortion and flat denominator below uses these original
weights, with no conditional renormalization of a cell.

For a pair $u<v$ of respective weights $p_u,p_v$, its exact Jensen curvature
kernel is
\begin{equation}
 K_{uv}(t)=\min\{p_u(t-u),p_v(v-t)\}\,
           \mathbf1_{[u,v]}(t).
 \label{eq:ext-pairkernel}
\end{equation}
For an arbitrary cell $S$, the kernel is
\[
 K_S(t)=\sum_{i\in S}p_i(x_i-t)_+-p_S(m_S-t)_+,
 \qquad D(S)=\int K_S(t)g_n(t)\rd t.
\]
These formulas follow by writing a twice differentiable convex generator,
up to an affine term, as an integral of $(x-t)_+$ against its curvature.

Set
\[
 P=D(AB)=D(CD),\quad q=D(BC),\quad
 T=D(ABC)=D(BCD),\quad J=D(AD).
\]
The equalities use reflection symmetry. Since
$128n^2b=128n^2d^2\to0$, equations
\eqref{eq:ext-tail}--\eqref{eq:ext-max} imply, eventually,
\begin{equation}
 I=[4b,d]\ne\varnothing,\qquad
 \int_I\kappa_n(t)\rd t
 \ge1-\tau_n-128n^2b\ge\tfrac12.
 \label{eq:ext-core}
\end{equation}
We may also assume $d\le1/16$ and $\tau_n\le d^2/2$.

The pair $AB$ has mean $4bd$. Thus on $I$, which lies above that mean,
$K_{AB}(t)=b(2d-t)\ge bd$. Conversely $K_{AB}\le2bd$ everywhere.
On $[0,2d]$, the reflected needle contributes at most $\tau_n$ to the
integrated curvature, since $1-2d\ge d$. Hence
\[
 \int_0^{2d}g_n\le1+\tau_n+2d^3\le2.
\]
Together with \eqref{eq:ext-core}, this proves
\begin{equation}
 \frac{bd}{2}\le P\le4bd.
 \label{eq:ext-nearpair}
\end{equation}
For the middle pair, $K_{BC}\le b/2$. The two needles contribute at most
$2\tau_n$ on $[2d,1-2d]$, and the constant curvature contributes at most
$d^2$. Therefore
\begin{equation}
 0<q\le\frac b2(2\tau_n+d^2)\le bd^2.
 \label{eq:ext-middlepair}
\end{equation}
The triple $ABC$ has mean $b/(1/2+b)\le2b$. For $t\in I$ this mean lies
below $t$, while both inner locations lie above $t$, so
\[
 K_{ABC}(t)=b(2d-t)+b(1-2d-t)=b(1-2t)\ge b/2.
\]
The outer pair has kernel $K_{AD}(t)=a\min(t,1-t)$. On $I$ this is at
least $4ab\ge b$, because $a\ge1/4$. Thus
\begin{equation}
 T\ge b/4,\qquad J\ge b/2.
 \label{eq:ext-coarse}
\end{equation}

We now cover all deterministic hierarchy choices, including noncontiguous
ones. A hierarchy with fewer than the allowed numbers of cells can be
refined to exactly three fine cells and two coarse cells without increasing
either cost: split fine cells within their current coarse cells, then split
the coarse level if necessary by grouping the three fine cells. Every exact
three-cell partition of four states merges one pair. Besides $BC$, each
pair has cost at least $P$. This is immediate for $AB,CD$; the pair $AC$
has the same endpoint weights as $AB$ and a more distant right endpoint,
so its kernel dominates $K_{AB}$ pointwise, and $BD$ is its reflection.
Finally, $J\ge b/2\ge P$ by \eqref{eq:ext-nearpair} and $d\le1/16$.

If the fine pair is not $BC$, then $\OPT_3\le q$ gives fine-level ratio
at least $P/q\ge1/(2d)$. If the fine pair is $BC$, its exact two-cell
coarsenings are $ABC\mid D$, $A\mid BCD$, and $AD\mid BC$. By
\eqref{eq:ext-coarse} each costs at least $b/4$. Since the available
coarse split $AB\mid CD$ gives $\OPT_2\le2P\le8bd$, this hierarchy has
coarse-level ratio at least $1/(32d)$. These two cases prove
\begin{equation}
 \rho_{\ph_n,\detp}\ge\frac1{32d_n},\qquad
 \rho_{\ph_n,\stp}\ge\frac1{64d_n},
 \label{eq:ext-lower}
\end{equation}
the second inequality following from Lemma~\ref{lem:ext-rounding}.
Notice that this proof needs neither exact flat optimizer identification nor
an assertion that contiguous deterministic hierarchies exhaust the optimum.

\begin{theorem}[Matching arbitrary-weight four-state lower bound]
\label{thm:ext-weightdegenerate}
Let $\widehat C_r^{\detp}$ and $\widehat C_r^{\stp}$ denote the two/three-output
four-state degree envelopes in which all positive source weights are allowed.
There is an absolute $c>0$ such that, for all sufficiently large $r$,
\[
 \widehat C_r^{\detp}\ge\widehat C_r^{\stp}
 \ge c\frac{r^2}{(\log r)^2}.
\]
Consequently no absolute-constant $O(r/\log r)$ bound holds in this weighted
four-state setting. The witnesses can have all posteriors in $[1/4,3/4]$
and polynomial generators inducing smooth strictly proper losses bounded
in $[0,1/2]$.
\end{theorem}
\begin{proof}
Equations \eqref{eq:ext-endpointneedle} and \eqref{eq:ext-lower} give the
claim at degrees $r=2n$, since $1/d_n=n^2/(16(\log n)^2)$.
For a degree-at-most-$r$ envelope take $n=\lfloor r/2\rfloor$.

To put all posteriors in the stated interior interval, use
$y=1/4+x/2$. The four locations become
$(1/4,1/4+d,3/4-d,3/4)$.
Let $M_n=\max_{t\in[-1/2,3/2]}g_n(t)$ and set
\[
 \widetilde\ph_n(y)=\frac{\ph_n(2y-1/2)}{4M_n}.
\]
The fact that $g_n$ is a sum of squares plus $d^2$ on the \emph{whole real
line} ensures $0<\widetilde\ph_n''\le1$ on $[0,1]$. This affine coordinate
change and positive rescaling preserve every price and curvature degree.
The generator remains symmetric about $1/2$. Put
$h(y)=\widetilde\ph_n(0)-\widetilde\ph_n(y)$ and define
\[
 \ell(y,1)=h(y)+(1-y)h'(y),\qquad
 \ell(y,0)=h(y)-yh'(y).
\]
Expected regret is $B_{\widetilde\ph_n}$, so the loss is strictly proper.
Symmetry and the supporting-line inequalities make both losses nonnegative;
their derivatives are $-(1-y)\widetilde\ph_n''(y)$ and
$y\widetilde\ph_n''(y)$, respectively. Their common maximum is
$\widetilde\ph_n'(1)=-\widetilde\ph_n'(0)
=\int_{1/2}^1\widetilde\ph_n''\le1/2$.
\end{proof}


The smallest mass in this family is $d_n^2=256(\log n)^4/n^4$.
The generators, weights, and posterior gaps all vary with $n$. This does
not determine a sharp joint minimum-mass/degree envelope. The normalization
preserves price, not absolute excess cost; explicit vanishing flat-cost
bounds appear after Theorem~\ref{thm:weighted-sharp}.
